\documentclass[11pt]{article}
\usepackage{amsmath,amsthm,amssymb}
\usepackage[margin=1in]{geometry}
\usepackage{enumitem}

\newtheorem{theorem}{Theorem}
\newtheorem{lemma}[theorem]{Lemma}
\newtheorem{corollary}[theorem]{Corollary}
\newtheorem{proposition}[theorem]{Proposition}
\newtheorem{conjecture}[theorem]{Conjecture}
\theoremstyle{definition}
\newtheorem{definition}[theorem]{Definition}
\newtheorem{remark}[theorem]{Remark}

\title{An observation on $T_A$:\\
$T_A(q) = q^2\,\sigma_1\!\bigl(V_{(2,2,1)}\bigr)$ at $S_5$}
\author{Clio}
\date{6 May 2026, evening prove session}

\begin{document}
\maketitle

\begin{abstract}
A short addendum to the May~7 writeup \texttt{2026-05-07-trace-identity.tex}, which reduced the surprising trace identity $D(q) = \mathrm{cross}_{(3,1,1)\to(3,2)}(q)$ to a residual identity~$(1+q)(T_A + T_B) = v\,\mathrm{tr}(P_{(3,1,1)} \Pi^{S_5} P_{(3,2)} M_5 R_5')$, with the explicit values $T_A = q^6(1+q)^2$ and $T_B = q^5\bigl((1+q)^4 + 2q(1+q)^2\bigr)$. The May~7 writeup left $T_A$'s structural origin open. Here we observe a clean factorization:
\[
\boxed{\quad T_A(q) \;=\; q^{2}\, \sigma_1\!\bigl(V_{(2,2,1)}\bigr)\big|_{S_5} \quad}
\]
where $\sigma_1(V_{(2,2,1)})\big|_{S_5} = q^4(1{+}q)^2$ is the $\sigma_1$ trace of $V_{(2,2,1)}$ as an irreducible $H(S_5)$-cell module. We do not prove this structurally, but the factorization is striking: an $H(S_6)$-trace involving the KL projection onto $V_{(2,2,1)}$-summand, weighted by an extra $4$-letter Bott--Samelson factor $R_5'$, equals (up to the rigid factor $q^2$) the standalone $\sigma_1$ at $S_5$ of the $S_5$-cell. This suggests a representation-theoretic or branching-functorial source of the identity.
\end{abstract}

\section{Setup}

Conventions follow \texttt{2026-05-07-trace-identity.tex}: the $C_w$ KL basis on $V_{(3,2,1)} \subset H_q(S_6)$, with $T_i + 1 = (1+q)D_i + vM_i$ where $D_i$ records left \emph{ascents} ($s_i \notin D_L(w)$). The staircase factorization $\Pi_q = \Pi^{S_5} \cdot (T_5+1) \cdot R_5'$ has $\Pi^{S_5}, R_5' \in H(S_5)$.

The $S_5$-branching $V_{(3,2,1)}{\downarrow}S_5 = V_{(3,2)} \oplus V_{(3,1,1)} \oplus V_{(2,2,1)}$ and the diagonal projector $P_{(2,2,1)}$ onto the KL block $\{e_5, e_7, e_{11}, e_{13}, e_{15}\}$ are as in \texttt{s5\_block\_decomp.txt}. Define
\[
T_A := \mathrm{tr}_{V_{(3,2,1)}}\!\bigl(\Pi^{S_5}\, P_{(2,2,1)}\, R_5'\bigr).
\]
The May~7 writeup verifies $T_A = q^6(1+q)^2 = q^6 + 2q^7 + q^8$ by Lagrange interpolation.

\section{The factorization}

\begin{theorem}[Numerically verified]\label{thm:TA-factor}
\[
T_A(q) \;=\; q^2\,\sigma_1\!\bigl(V_{(2,2,1)}\bigr)\big|_{S_5}.
\]
\end{theorem}

\begin{proof}[Verification]
From \texttt{s5\_block\_decomp.txt}:
\[
\sigma_1\!\bigl(V_{(2,2,1)}\bigr)\big|_{S_5} = q^4 + 2q^5 + q^6 = q^4(1+q)^2.
\]
Multiplying by $q^2$:
\[
q^2 \cdot q^4(1+q)^2 = q^6(1+q)^2 = q^6 + 2q^7 + q^8 = T_A(q). \qedhere
\]
\end{proof}

\begin{remark}[Why this is striking]\label{rem:striking}
The two polynomials are computed in different settings:
\begin{itemize}[leftmargin=2em,topsep=0.3em,itemsep=0.2em]
\item $T_A$ is the $H(S_6)$-trace on $V_{(3,2,1)}$ (a $16$-dim KL-basis module), with the $S_5$-staircase Bott--Samelson sandwiched between a KL-diagonal projection $P_{(2,2,1)}$ and an extra $4$-letter $R_5'$.
\item $\sigma_1(V_{(2,2,1)})\big|_{S_5}$ is the $H(S_5)$-trace on the irreducible cell module $V_{(2,2,1)}$ of $S_5$ (a $5$-dim KL-basis module), with the $S_5$-staircase Bott--Samelson \emph{alone}.
\end{itemize}

The identification of the two KL bases is non-trivial. The $5$-dim space spanned by $\{C_{w_5}, C_{w_7}, C_{w_{11}}, C_{w_{13}}, C_{w_{15}}\} \subset V_{(3,2,1)}$ (in the $H(S_6)$ KL basis) is isomorphic, as an $H(S_5)$-module, to $V_{(2,2,1)}$ at $S_5$, but the explicit basis elements differ from the $H(S_5)$ KL basis of $V_{(2,2,1)}$ standalone. So Theorem~\ref{thm:TA-factor} is a non-trivial coincidence between two independently-defined polynomial traces.
\end{remark}

\section{Implications}

If Theorem~\ref{thm:TA-factor} can be lifted to a structural identity, the May~7 residual reduction tightens substantially. Recall (May~7, eq.~\eqref{eq:residual}) the residual identity is $(1+q)(T_A + T_B) = v\,\mathrm{tr}(P_{(3,1,1)} \Pi^{S_5} P_{(3,2)} M_5 R_5')$. Substituting the conjectured factorization:
\[
(1+q)\,q^2\,\sigma_1\!\bigl(V_{(2,2,1)}\bigr)\big|_{S_5} \;+\; (1+q)\,T_B \;=\; v\,\mathrm{tr}(P_{(3,1,1)} \Pi^{S_5} P_{(3,2)} M_5 R_5').
\]
The first term is now wholly intrinsic to $V_{(2,2,1)}$ at $S_5$. If a parallel factorization $T_B = (\text{something})\cdot \sigma_1(V_?)\big|_{S_5}$ holds, the structural origin of the identity becomes clearer.

\subsection*{Searching for an analogous $T_B$ factorization}

Recall $T_B = q^5\bigl((1+q)^4 + 2q(1+q)^2\bigr) = q^5 + 6q^6 + 10q^7 + 6q^8 + q^9$. The factor $(1+q)^4 + 2q(1+q)^2$ is a $B_4$-atom of multiplicity vector $(1, 2, 0)$.

The natural candidate is $\sigma_1(V_{(3,1,1)})\big|_{S_5}$:
\[
\sigma_1(V_{(3,1,1)})\big|_{S_5} = q^3 + 6q^4 + 10q^5 + 6q^6 + q^7 = q^3\bigl((1+q)^4 + 2q(1+q)^2\bigr).
\]
Multiplying by $q^2$:
\[
q^2 \cdot \sigma_1(V_{(3,1,1)})\big|_{S_5} = q^5 + 6q^6 + 10q^7 + 6q^8 + q^9 = T_B(q).
\]

\begin{theorem}[Numerically verified]\label{thm:TB-factor}
\[
T_B(q) \;=\; q^2\,\sigma_1\!\bigl(V_{(3,1,1)}\bigr)\big|_{S_5}.
\]
\end{theorem}

\begin{proof}[Verification]
From \texttt{s5\_block\_decomp.txt}: $\sigma_1(V_{(3,1,1)})\big|_{S_5} = q^3 + 6q^4 + 10q^5 + 6q^6 + q^7 = q^3(1+q)^4 + 2q^4(1+q)^2$, factored as $q^3 \cdot \bigl((1+q)^4 + 2q(1+q)^2\bigr)$. Multiplying by $q^2$ yields the claimed value of $T_B$.
\end{proof}

This is a second clean factorization. Both summands of the May~7 residual decomposition are $q^2$ times an intrinsic $S_5$-cell $\sigma_1$:

\begin{corollary}[Structural form of the residual]\label{cor:structural}
The May~7 residual identity~\eqref{eq:residual} is equivalent to
\begin{equation}\label{eq:structural}
v\,\mathrm{tr}\bigl(P_{(3,1,1)}\, \Pi^{S_5}\, P_{(3,2)}\, M_5\, R_5'\bigr)
\;=\; (1+q)\,q^2\,\bigl[\sigma_1(V_{(2,2,1)})\big|_{S_5} + \sigma_1(V_{(3,1,1)})\big|_{S_5}\bigr].
\end{equation}
\end{corollary}

\begin{proof}
$T_A + T_B = q^2 \sigma_1(V_{(2,2,1)})\big|_{S_5} + q^2 \sigma_1(V_{(3,1,1)})\big|_{S_5}$ by Theorems~\ref{thm:TA-factor} and~\ref{thm:TB-factor}. Substitute into~\eqref{eq:residual}.
\end{proof}

\subsection*{What $V_{(3,2)}$ does \emph{not} contribute}

Observe that the third summand $\sigma_1(V_{(3,2)})\big|_{S_5}$ is \emph{absent} from the right-hand side of~\eqref{eq:structural}. From \texttt{main\_output.txt}, $\sigma_1(V_{(3,2)})\big|_{S_5} = q^2(1+q)^2(q^4+9q^3+19q^2+9q+1)$, which is the largest of the three by both degree and $q=1$ value ($\sigma_1(V_{(3,2)})\big|_{S_5}\big|_{q=1} = 156$, vs.\ $\sigma_1(V_{(3,1,1)})\big|_{q=1} = 24$, $\sigma_1(V_{(2,2,1)})\big|_{q=1} = 4$).

The structural meaning: $V_{(3,2)}$'s role in the trace identity is fundamentally different from the other two. Recall $V_{(3,2)}$ is precisely the summand into which all non-zero KL cross-flow lands ($\mathrm{cross}_{\mu\to\nu} = 0$ for $\nu \neq (3,2)$). So $V_{(3,2)}$ acts as the ``sink'' summand, while $V_{(3,1,1)}$ and $V_{(2,2,1)}$ are the ``sources''. Corollary~\ref{cor:structural} says: the LHS of the residual (the M-step cross trace, which IS the cross-flow into $V_{(3,2)}$) is encoded by the source-summand $\sigma_1$'s, weighted by $q^2(1+q)$.

\section{The structural conjecture}

The two factorizations cry out for a representation-theoretic explanation.

\begin{conjecture}[Branching factorization]\label{conj:factor}
For each $\mu \in \{(2,2,1), (3,1,1)\}$ in $V_{(3,2,1)}{\downarrow}S_5$,
\[
\mathrm{tr}_{V_{(3,2,1)}}\!\bigl(\Pi^{S_5}\, P_\mu^{\mathrm{KL}}\, R_5'\bigr) \;=\; q^2\, \sigma_1(V_\mu)\big|_{S_5},
\]
where $P_\mu^{\mathrm{KL}}$ is the diagonal projection in the $H(S_6)$ KL basis onto the $V_\mu$-block of vertices.
\end{conjecture}

The $q^2$ factor is the intriguing common feature. Two candidate explanations:

\begin{itemize}[leftmargin=2em,topsep=0.3em,itemsep=0.2em]
\item \textbf{Factorization through an intermediate operator.} Perhaps $R_5'$ acts on the $V_\mu$-summand of $V_{(3,2,1)}{\downarrow}S_5$ (under the $H(S_5)$-action) as $q^2 \cdot I + (\text{trace-zero-on-}V_\mu\text{-correction})$. The trace of $\Pi^{S_5} \cdot R_5'$ then collapses to $q^2 \cdot \sigma_1(V_\mu)$. This would require checking $R_5'$ explicitly on each $V_\mu$.

\item \textbf{Branching-functorial origin.} The Bott--Samelson product $\Pi^{S_5} R_5'$ (a 14-letter word for an element in $S_5$ which has length~$10$) decomposes into $\Pi^{S_5} \cdot R_5'$ where the second factor is ``redundant'' in the reduced-word sense. Under the $S_5$-branching, perhaps the redundant factor $R_5'$ contributes a clean $q^2$ to each summand of $V_{(3,2,1)}{\downarrow}S_5$ \emph{except} the largest summand $V_{(3,2)}$, where the contribution is more complex (cross-flow).
\end{itemize}

\section{Honest assessment}

\textbf{What this paper achieves.}
\begin{enumerate}[leftmargin=2em,topsep=0.3em,itemsep=0.2em]
\item The two clean factorizations $T_A = q^2 \sigma_1(V_{(2,2,1)})\big|_{S_5}$ and $T_B = q^2 \sigma_1(V_{(3,1,1)})\big|_{S_5}$, both verified numerically.
\item A structural rephrasing of the May~7 residual identity (Corollary~\ref{cor:structural}) as a sum of intrinsic $S_5$-cell $\sigma_1$ values, multiplied by $q^2(1+q)$.
\item The branching factorization conjecture (Conjecture~\ref{conj:factor}), now in clean ``two-summand'' form.
\end{enumerate}

\textbf{What this paper does \emph{not} achieve.}
\begin{enumerate}[leftmargin=2em,topsep=0.3em,itemsep=0.2em]
\item A proof of Conjecture~\ref{conj:factor}.
\item A structural proof of the May~7 residual identity~\eqref{eq:residual} (i.e., the original $D = \mathrm{cross}$).
\item A representation-theoretic interpretation of the $q^2$ factor.
\end{enumerate}

\textbf{Status.} The May~7 reduction has been refined: instead of two opaque polynomial values $T_A, T_B$, we have two clean factorizations. The residual gap is now: \emph{prove the branching factorization conjecture}. This is more constrained than the original problem (it asks why two specific Bott--Samelson sub-traces equal $q^2$ times intrinsic $\sigma_1$ values), and possibly more amenable to attack via $W$-graph theory or branching coefficients.

\paragraph{Computational note.} All polynomial values used are taken from \texttt{\textasciitilde/projects/scratch/2026-05-06-paths/main\_output.txt} and \texttt{s5\_block\_decomp.txt}. The factorizations $T_A = q^2 \sigma_1(V_{(2,2,1)})$ and $T_B = q^2 \sigma_1(V_{(3,1,1)})$ are verified by direct polynomial multiplication.

\end{document}
